\section{Conclusion}
\label{sec:conclusion}

This paper investigated secure UAV communication with digital-twin uncertainty and limited synchronization resources. We formulated a joint synchronization and control problem in which two cooperative UAVs select synchronization bandwidth, motion, source power, and jamming power while serving ground users in the presence of a mobile eavesdropper. We proposed an empirical-risk-aware joint rollout controller and validated a holdout-fitted secrecy-loss risk estimator.

The final results support the value of joint synchronization-control design. At the calibrated non-saturated stress threshold, the proposed controller improves average secrecy rate over periodic synchronization by 0.2683 and reduces secrecy outage by 0.2141. A stress-threshold sweep shows that the secrecy-rate gain is stable across the tested $R_{\min}$ values, while outage reduction is most meaningful within the informative non-saturated threshold range. Additional stress-scenario experiments show that it outperforms SCA baselines, no-twin rollout, no-synchronization rollout, and fixed-period rollout, while remaining below an oracle-state upper reference. A lightweight PPO diagnostic run did not learn a competitive policy within 200 episodes, so stronger DRL comparisons remain future work. The risk-estimator validation results show empirical cover rates above the diagnostic 0.95 target on the evaluated holdout groups.

The method is not yet a low-latency onboard solution because its rollout search is computationally expensive. The risk estimator is also empirical and should be recalibrated under major distribution shifts. Even with these limitations, the study provides evidence that digital-twin synchronization should be treated as an active security-control variable rather than as a fixed background maintenance process.
